% Source: nn_activation_functions.tex
\begin{notechapter}
\chapter{A Catalogue of Activation Functions}\label{ch:activation-catalogue}
\chaptersubtitle{Advantages and disadvantages}

\section{Introduction}
Activation functions are a critical component of neural networks. They introduce
non-linearity, enabling the network to learn complex patterns. This document
provides an overview of the most well-known activation functions, their
mathematical definitions, advantages, disadvantages, and typical use cases. It
also includes a detailed explanation of the vanishing gradient problem,
saturation, and the issue of non-zero-centered activations.

\section{Sigmoid}
The sigmoid function is defined as:
\[
\sigma(z) = \frac{1}{1 + e^{-z}}.
\]
Its derivative is:
\[
\sigma'(z) = \sigma(z) (1 - \sigma(z)).
\]

\subsection{Advantages}
\begin{itemize}[leftmargin=*]
    \item \textbf{Smooth Gradient}: The sigmoid function is differentiable,
    providing a smooth gradient for backpropagation.
    \item \textbf{Output Range}: The output is bounded between 0 and 1, making it
    suitable for binary classification tasks.
\end{itemize}

\subsection{Disadvantages}
\begin{itemize}[leftmargin=*]
    \item \textbf{Vanishing Gradient}: For very high or low inputs, the gradient
    becomes very small, slowing down learning.
    \item \textbf{Saturation}: The sigmoid function saturates for extreme inputs,
    causing the gradient to approach 0.
    \item \textbf{Non-Zero-Centered}: The output is always positive, which can
    lead to inefficient weight updates (explained in detail below).
    \item \textbf{Computationally Expensive}: The exponential function is
    relatively slow to compute.
\end{itemize}

\subsection{Typical Use Cases}
\begin{itemize}[leftmargin=*]
    \item Binary classification (output layer).
    \item Historical use in early neural networks.
\end{itemize}

\section{Hyperbolic Tangent (Tanh)}
The tanh function is defined as:
\[
\tanh(z) = \frac{e^z - e^{-z}}{e^z + e^{-z}}.
\]
Its derivative is:
\[
\tanh'(z) = 1 - \tanh^2(z).
\]

\subsection{Advantages}
\begin{itemize}[leftmargin=*]
    \item \textbf{Zero-Centered}: The output is centered around 0, leading to
    more efficient weight updates.
    \item \textbf{Stronger Gradient}: The gradient is steeper than sigmoid,
    reducing the vanishing gradient problem.
\end{itemize}

\subsection{Disadvantages}
\begin{itemize}[leftmargin=*]
    \item \textbf{Vanishing Gradient}: Like sigmoid, tanh suffers from vanishing
    gradients for extreme inputs.
    \item \textbf{Saturation}: The tanh function saturates for extreme inputs,
    causing the gradient to approach 0.
    \item \textbf{Computationally Expensive}: The exponential function is slow
    to compute.
\end{itemize}

\subsection{Typical Use Cases}
\begin{itemize}[leftmargin=*]
    \item Hidden layers in neural networks.
    \item Tasks requiring zero-centered outputs.
\end{itemize}

\section{Rectified Linear Unit (ReLU)}
The ReLU function is defined as:
\[
\text{ReLU}(z) = \max(0, z).
\]
Its derivative is:
\[
\text{ReLU}'(z) = \begin{cases}
    1 & \text{if } z > 0, \\
    0 & \text{if } z \leq 0.
\end{cases}
\]

\subsection{Advantages}
\begin{itemize}[leftmargin=*]
    \item \textbf{Computationally Efficient}: ReLU is simple and fast to compute.
    \item \textbf{Sparsity}: ReLU outputs exactly 0 for negative inputs, leading
    to sparse activations and improved efficiency.
    \item \textbf{Reduced Vanishing Gradient}: The gradient is 1 for positive
    inputs, mitigating the vanishing gradient problem.
\end{itemize}

\subsection{Disadvantages}
\begin{itemize}[leftmargin=*]
    \item \textbf{Dying ReLU Problem}: Neurons can get stuck outputting 0 and
    stop learning (dead neurons).
    \item \textbf{Non-Zero-Centered}: Like sigmoid, ReLU outputs are always
    non-negative, which can lead to inefficient weight updates (explained below).
\end{itemize}

\subsection{Typical Use Cases}
\begin{itemize}[leftmargin=*]
    \item Hidden layers in deep neural networks.
    \item Tasks requiring fast computation and sparse activations.
\end{itemize}

\section{Why Non-Zero-Centered Activations Are a Disadvantage}
Non-zero-centered activation functions (e.g., sigmoid, ReLU) can lead to
inefficient weight updates during training. This is because the gradient of the
loss with respect to the weights (\(\frac{\partial L}{\partial W_{ij}}\)) depends
on the activations of the previous layer (\(a_i\)).

\subsection{Mathematical Explanation}
The gradient of the loss \(L\) with respect to a weight \(W_{ij}\) is given by:
\[
\frac{\partial L}{\partial W_{ij}} = \frac{\partial L}{\partial z_j} \cdot a_i,
\]
where:
\begin{itemize}[leftmargin=*]
    \item \(\frac{\partial L}{\partial z_j}\) is the gradient of the loss with
    respect to the pre-activation \(z_j\).
    \item \(a_i\) is the activation of the \(i\)-th neuron in the previous layer.
\end{itemize}

If \(a_i\) is always non-negative (e.g., sigmoid or ReLU), the sign of
\(\frac{\partial L}{\partial W_{ij}}\) depends only on \(\frac{\partial L}{\partial z_j}\).
This can lead to \textbf{inefficient weight updates} because all gradients for the
weights will have the \textbf{same sign}, causing the optimization process to
zig-zag and slow down convergence.

\subsection{Example: ReLU vs. Tanh}
Consider a simple neural network with one hidden layer and ReLU activation:
\begin{itemize}[leftmargin=*]
    \item If the input to the ReLU is positive, the activation is positive.
    \item If the input is negative, the activation is 0.
\end{itemize}

This means:
\begin{itemize}[leftmargin=*]
    \item The gradients for the weights will always have the same sign (positive
    or negative), leading to inefficient updates.
\end{itemize}

In contrast, with tanh:
\begin{itemize}[leftmargin=*]
    \item The activations can be both positive and negative.
    \item The gradients for the weights can have both signs, leading to more
    balanced updates.
\end{itemize}

\subsection{Impact on Training}
\begin{itemize}[leftmargin=*]
    \item \textbf{Inefficient Weight Updates}: Non-zero-centered activations
    cause all gradients to have the same sign, leading to zig-zagging updates
    and slower convergence.
    \item \textbf{Zero-Centered Activations}: Zero-centered activations (e.g.,
    tanh) allow gradients to have both positive and negative signs, leading to
    more balanced updates and faster convergence.
\end{itemize}

\subsection{Conclusion}
Non-zero-centered activations can slow down training due to inefficient weight
updates. Zero-centered activations (e.g., tanh) or alternatives like Leaky ReLU
and ELU can mitigate this issue by allowing both positive and negative outputs,
leading to faster and more stable convergence.

\section{Vanishing Gradient Problem and Saturation}
The vanishing gradient problem and saturation are two related but distinct issues
that can arise in neural networks, particularly when using activation functions
like sigmoid and tanh.

\subsection{Vanishing Gradient Problem}
The vanishing gradient problem occurs when the gradients of the loss with respect
to the weights become extremely small during backpropagation. This happens because
the gradients are computed using the chain rule, which involves multiplying many
derivatives together. If these derivatives are small, their product can become
vanishingly small, effectively stopping the weights from updating.

\subsubsection{Causes}
\begin{itemize}[leftmargin=*]
    \item \textbf{Activation Functions}: Activation functions like sigmoid and tanh
    have derivatives that are close to 0 for large positive or negative inputs.
    \item \textbf{Deep Networks}: In deep networks, the gradients are multiplied
    through many layers. If each layer's gradient is small, the overall gradient
    can vanish.
\end{itemize}

\subsubsection{Impact}
\begin{itemize}[leftmargin=*]
    \item The weights are updated very slowly or not at all, causing the network
    to stop learning.
    \item This is particularly problematic in deep networks, where gradients must
    propagate through many layers.
\end{itemize}

\subsection{Saturation}
Saturation refers to the phenomenon where the output of an activation function
becomes insensitive to changes in the input. This happens when the input to the
activation function is very large or very small, causing the function to "flatline."

\subsubsection{Causes}
\begin{itemize}[leftmargin=*]
    \item \textbf{Activation Functions}: Sigmoid and tanh are prone to saturation
    because their outputs approach fixed values (0 or 1 for sigmoid, -1 or 1 for
    tanh) for extreme inputs.
    \item \textbf{Weight Initialization}: Poor weight initialization can cause
    neurons to start in the saturated region, making it difficult for them to
    recover.
\end{itemize}

\subsubsection{Impact}
\begin{itemize}[leftmargin=*]
    \item Saturated neurons have gradients close to 0, which contributes to the
    vanishing gradient problem.
    \item The network becomes less expressive because saturated neurons stop
    responding to changes in the input.
\end{itemize}

\subsection{Relationship Between Vanishing Gradients and Saturation}
\begin{itemize}[leftmargin=*]
    \item \textbf{Saturation Leads to Vanishing Gradients}: When neurons saturate,
    their gradients become very small. During backpropagation, these small gradients
    are multiplied through many layers, leading to the vanishing gradient problem.
    \item \textbf{Vanishing Gradients Exacerbate Saturation}: If the gradients vanish,
    the weights are not updated effectively, and neurons may remain in the saturated
    region.
\end{itemize}

\subsection{Example: Sigmoid Activation}
Consider the sigmoid activation function:
\[
\sigma(z) = \frac{1}{1 + e^{-z}}.
\]
Its derivative is:
\[
\sigma'(z) = \sigma(z) (1 - \sigma(z)).
\]
\begin{itemize}[leftmargin=*]
    \item For large positive \(z\), \(\sigma(z) \approx 1\) and \(\sigma'(z) \approx 0\).
    \item For large negative \(z\), \(\sigma(z) \approx 0\) and \(\sigma'(z) \approx 0\).
\end{itemize}
In both cases, the neuron is saturated, and the gradient is close to 0, contributing
to the vanishing gradient problem.

\subsection{Independent Vanishing Gradient Problem}
Even if the activation function is not fully saturated, the vanishing gradient problem
can still occur due to the multiplication of small gradients. For example:
\begin{itemize}[leftmargin=*]
    \item The maximum gradient of the sigmoid function is 0.25.
    \item In a deep network, multiplying many such small gradients can cause the
    overall gradient to vanish.
\end{itemize}

\subsection{Mitigating Vanishing Gradients and Saturation}
\begin{itemize}[leftmargin=*]
    \item \textbf{Activation Functions}: Use activation functions that are less prone
    to saturation and vanishing gradients, such as ReLU, Leaky ReLU, or ELU.
    \item \textbf{Weight Initialization}: Use proper weight initialization techniques
    to avoid starting in the saturated region.
    \item \textbf{Batch Normalization}: Normalize the inputs to each layer to prevent
    saturation and improve gradient flow.
    \item \textbf{Gradient Clipping}: Clip gradients to prevent them from becoming
    too small or too large.
\end{itemize}

\section{Conclusion}
Activation functions play a crucial role in neural networks, introducing
non-linearity and enabling complex learning. Each activation function has its
advantages and disadvantages, making it suitable for specific tasks. By
understanding these trade-offs, you can choose the best activation function for
your neural network.
\end{notechapter}
